\documentclass[11pt]{article}
\usepackage[margin=1in]{geometry}
\usepackage{amsmath,amssymb,amsthm,mathtools}
\usepackage{graphicx}
\usepackage{xurl}
\usepackage[colorlinks=true,linkcolor=blue,citecolor=blue,urlcolor=blue]{hyperref}
\newtheorem{theorem}{Theorem}[section]
\newtheorem{lemma}[theorem]{Lemma}
\newtheorem{corollary}[theorem]{Corollary}
\newtheorem{remark}[theorem]{Remark}
\newcommand{\A}{\mathcal A}
\newcommand{\Ah}{\mathcal A^{\#}}
\newcommand{\tp}{\widehat\otimes}
\newcommand{\e}{\mathbf e}
\emergencystretch=1em
\title{A Two-Way Link Between Pseudo-Amenability\\
and Approximate Semi-Amenability}
\author{Automated mathematics run \texttt{fa\_banach\_001}, lane 3}
\date{11 August 2026}

\begin{document}
\maketitle

\begin{abstract}
Ghahramani and Loy ask for the connection between pseudo-amenability and
approximate semi-amenability.  We prove that every pseudo-amenable Banach
algebra is approximately semi-amenable.  Conversely, an approximately
semi-amenable Banach algebra with a central approximate identity is
pseudo-amenable; the approximate identity need not be bounded.  It follows
that the two notions are equivalent for every commutative Banach algebra.
The unrestricted noncommutative converse is not settled here.  Status:
\emph{candidate partial result, likely valid, pending human review}.
\end{abstract}

\section{Source question and definitions}
The source is F.~Ghahramani and R.~J.~Loy,
\emph{Approximate semi-amenability of Banach algebras}, arXiv:1910.03775
\cite{GL}.  Its second open problem asks the question shown in
Figure~\ref{fig:source}.

\begin{figure}[ht]
 \centering
 \includegraphics[width=0.94\textwidth]{figures/open_problem_crop.png}
 \caption{The exact source question, cropped from arXiv PDF page~24.}
 \label{fig:source}
\end{figure}

Write $\Ah$ for the unitization of a Banach algebra $\A$, $\e$ for the
adjoined identity, $\Ah\tp\Ah$ for the projective tensor product, and
$\pi(x\otimes y)=xy$.  The usual module actions are
\[
 a\cdot(x\otimes y)=ax\otimes y,
 \qquad (x\otimes y)\cdot a=x\otimes ya.
\]
The algebra $\A$ is \emph{pseudo-amenable} if it has an approximate diagonal:
a net $(m_\lambda)\subset \A\tp\A$ such that, for every $a\in\A$,
\begin{equation}\label{eq:pseudo}
 a\cdot m_\lambda-m_\lambda\cdot a\longrightarrow0,
 \qquad \pi(m_\lambda)a\longrightarrow a.
\end{equation}

Theorems 2.6 and 2.11 and Remark 2.7 of the source give the criterion we use:
$\A$ is approximately semi-amenable if and only if there are nets
$\xi_i,\eta_i\in\Ah\tp\Ah$ such that
\begin{equation}\label{eq:semi}
 a\cdot\xi_i-\eta_i\cdot a\longrightarrow0\quad(a\in\A),
 \qquad \pi(\xi_i)=\pi(\eta_i)=\e.
\end{equation}
The passage from approximate semi-amenability to tensors in the actual
projective tensor product uses the source's equivalence with approximate
semi-contractibility.

\section{Proof intuition}
An approximate diagonal is one tensor that is almost central.  The source
criterion allows two tensors: left multiplication acts on $\xi_i$, while
right multiplication acts on $\eta_i$.  Starting with one approximate
diagonal, there is a direct unitization trick that manufactures this pair,
which proves the unconditional forward implication.

For the reverse direction, the mixed relation in \eqref{eq:semi} implies that
$\xi_i-\eta_i$ disappears after multiplication by fixed algebra elements on
both sides.  A central approximate identity can therefore localize $\xi_i$
from $\Ah\tp\Ah$ into $\A\tp\A$ without creating commutator terms.  The
localizers may be unbounded because they are chosen first and the tensor
index is chosen afterward.  This order of choices is the key point.

\section{The unconditional implication}
\begin{theorem}\label{thm:forward}
Every pseudo-amenable Banach algebra is approximately semi-amenable.
\end{theorem}
\begin{proof}
Let $(m_\lambda)\subset\A\tp\A$ satisfy \eqref{eq:pseudo}, and put
$u_\lambda=\pi(m_\lambda)$.  Applying $\pi$ to the first limit in
\eqref{eq:pseudo} gives
\[
 au_\lambda-u_\lambda a\longrightarrow0.
\]
Since $u_\lambda a\to a$, it follows that $au_\lambda\to a$ as well.  In
$\Ah\tp\Ah$, define
\begin{align*}
 \xi_\lambda&=\e\otimes\e-u_\lambda\otimes\e+m_\lambda,\\
 \eta_\lambda&=\e\otimes\e-\e\otimes u_\lambda+m_\lambda.
\end{align*}
Their products are exactly $\e$.  For $a\in\A$, direct expansion gives
\begin{align*}
 a\cdot\xi_\lambda-\eta_\lambda\cdot a
  ={}&(a-au_\lambda)\otimes\e
      +\e\otimes(u_\lambda a-a)\\
    &+(a\cdot m_\lambda-m_\lambda\cdot a),
\end{align*}
which tends to zero.  Criterion \eqref{eq:semi} completes the proof.
\end{proof}

\section{A converse with an unbounded central approximate identity}
We first isolate the cancellation hidden in the two-tensor relation.

\begin{lemma}\label{lem:annihilation}
Suppose $\xi_i,\eta_i$ obey the first limit in \eqref{eq:semi}, and set
$r_i(a)=a\cdot\xi_i-\eta_i\cdot a$.  Then for fixed $b,c\in\A$,
\begin{equation}\label{eq:annihilation}
 b\cdot(\xi_i-\eta_i)\cdot c
 =r_i(b)\cdot c-r_i(bc)+b\cdot r_i(c)\longrightarrow0.
\end{equation}
\end{lemma}
\begin{proof}
The identity follows by expanding its right-hand side; each of the three
terms tends to zero after multiplication by the fixed elements $b,c$.
\end{proof}

\begin{theorem}\label{thm:converse}
Let $\A$ have a central approximate identity, not assumed bounded.  If $\A$
is approximately semi-amenable, then $\A$ is pseudo-amenable.
\end{theorem}
\begin{proof}
Take $\xi_i,\eta_i$ satisfying \eqref{eq:semi}.  We construct approximate
diagonals locally.  Fix a finite set $F\subset\A$ and $\varepsilon>0$.
Choose a central approximate-identity element $f\in\A$ so that $fa$ is close
to $a$ for every $a\in F$.  Then choose another central
approximate-identity element $d\in\A$ so that $d(fa)$ is close to $fa$ for
every $a\in F$.  The choices can be made so that
\begin{equation}\label{eq:product-ai}
 \|dfa-a\|<\varepsilon\qquad(a\in F).
\end{equation}
No common norm bound for $d$ and $f$ is needed.

After $d$ and $f$ are fixed, choose $i$ sufficiently late and put
\[
 M=d\cdot\xi_i\cdot f.
\]
Multiplication by $d$ on the first tensor factor and by $f$ on the second
shows that $M\in\A\tp\A$, even though $\xi_i\in\Ah\tp\Ah$.  Moreover,
\begin{equation}\label{eq:product}
 \pi(M)=d\pi(\xi_i)f=df.
\end{equation}
Since $d$ and $f$ are central, for $a\in F$ we have
\begin{align}
 a\cdot M-M\cdot a
 &=d\cdot(a\cdot\xi_i-\xi_i\cdot a)\cdot f\notag\\
 &=d\cdot r_i(a)\cdot f
   +d\cdot(\eta_i-\xi_i)\cdot af.\label{eq:commutator}
\end{align}
The first term tends to zero by \eqref{eq:semi}; the second tends to zero by
Lemma~\ref{lem:annihilation}, applied with the now-fixed elements $d$ and
$af$.  We can therefore choose $i$ so that
$\|a\cdot M-M\cdot a\|<\varepsilon$ for every $a\in F$.  Equations
\eqref{eq:product-ai} and \eqref{eq:product} also give
$\|\pi(M)a-a\|<\varepsilon$.

Direct the triples consisting of a finite subset of $\A$, a positive
tolerance, and the local choices above.  The resulting net $(M)$ in
$\A\tp\A$ satisfies \eqref{eq:pseudo}; hence $\A$ is pseudo-amenable.
\end{proof}

\begin{corollary}\label{cor:commutative}
For every commutative Banach algebra, approximate semi-amenability is
equivalent to pseudo-amenability.
\end{corollary}
\begin{proof}
The forward direction from pseudo-amenability is
Theorem~\ref{thm:forward}.  Conversely, Theorem 2.13 of the source says that
an approximately semi-amenable algebra has left and right approximate
identities.  In a commutative algebra, either is a central two-sided
approximate identity, so Theorem~\ref{thm:converse} applies.
\end{proof}

\begin{remark}
The unboundedness allowance is substantive.  The source proves that a
central \emph{bounded} approximate identity upgrades approximate
semi-amenability to approximate amenability.  Theorem~\ref{thm:converse}
requires only an unbounded central approximate identity and reaches the
weaker, natural conclusion of pseudo-amenability.
\end{remark}

\begingroup
\small
\section{Verification, upgrade attempts, and scope}
Every displayed tensor identity was expanded independently.  In particular,
\eqref{eq:annihilation} is exact, $M$ really belongs to $\A\tp\A$, and the
two central approximate-identity elements are fixed before the tensor index;
this prevents their potentially large norms from spoiling convergence.

Four focused upgrade attempts were made.  First, localizing by the separate
left and right approximate identities available in every approximately
semi-amenable algebra creates uncontrolled commutators when the localizers
are noncentral.  Second, decomposing $\xi_i,\eta_i$ into their four unitization
components leaves asymmetric correction terms; flip symmetrization cancels
them only in the commutative case.  Third, the standard finite-corner
diagonals for Schatten classes have growing projective norm, so arbitrary
Schatten tails do not support the rate needed to pass to the limit.  Fourth,
known examples separating pseudo-amenability from approximate amenability do
not separate the two notions studied here, because
Theorem~\ref{thm:forward} puts them on the approximately semi-amenable side.
The unrestricted noncommutative converse and the source's Schatten-class
question therefore remain open in this packet.

A bounded novelty search on 11 August 2026 used the exact arXiv id, title,
open-question wording, and combinations of ``pseudo-amenable'',
``approximately semi-amenable'', ``central approximate identity'', and
``commutative Banach algebra''.  It found the source and general
pseudo-amenability literature, including \cite{GZ}, but no direct statement
of Theorems~\ref{thm:forward} or \ref{thm:converse}, or of
Corollary~\ref{cor:commutative}.  This was not an exhaustive literature
review.

Human review should focus on the local-to-net construction in
Theorem~\ref{thm:converse} and on whether the three statements already occur
under different terminology in later Banach-algebra literature.

\begin{thebibliography}{9}
\bibitem{GL}
F.~Ghahramani and R.~J.~Loy,
\emph{Approximate semi-amenability of Banach algebras},
Semigroup Forum \textbf{101} (2020), 358--384,
\href{https://doi.org/10.1007/s00233-020-10094-6}{doi:10.1007/s00233-020-10094-6};
arXiv:1910.03775.

\bibitem{GZ}
F.~Ghahramani and Y.~Zhang,
\emph{Pseudo-amenable and pseudo-contractible Banach algebras},
Math. Proc. Cambridge Philos. Soc. \textbf{142} (2007), 111--123,
\href{https://doi.org/10.1017/S0305004106009649}{doi:10.1017/S0305004106009649}.
\end{thebibliography}
\endgroup

\end{document}
